% STATUS: DRAFT
\chapter{Gauge theory warm-up: edge modes and centres}\label{ch:gauge_warmup}

\begin{physmotiv}
Before gravity, consider the simplest place where ``constraints change the algebra of a region'': ordinary gauge theory. Gauss's law forces gauge-invariant operators in a region to commute with the electric flux through its boundary. This flux is then a \emph{central} element: a classical variable attached to the boundary. The Hilbert space does not factorize, the entanglement entropy picks up a classical Shannon term, and one is led to \emph{edge modes}. The lesson, that constraints decide which operators are observable in a region and can produce a centre, is the one that gravity will teach again, with the observer playing the role of the boundary.
\end{physmotiv}

\section{A lattice $U(1)$ example}

Take $U(1)$ gauge theory on a lattice. On each link $\ell$ there is a quantum rotor with electric field $E_\ell\in\Z$ and a unitary $U_\ell=\ee^{\ii A_\ell}$ raising it. Gauss's law at each vertex $v$ reads $\sum_{\ell\ni v}\pm E_\ell=\rho_v$ (charge). Gauge-invariant operators are generated by $E_\ell$, by Wilson loops $\prod U_\ell$ around closed paths, and by open Wilson lines ending on charges.

Cut the lattice into regions $A$ and $B$ along a set of links $\Sigma$. Gauge-invariant operators supported in $A$ include everything built from links in $A$; those on $\Sigma$ are shared. A Wilson line from $A$ to $B$ crossing the cut is \emph{not} in either algebra, yet cannot be dropped. On the other hand, the electric field $E_\ell$ on a cut link commutes with all gauge-invariant operators in $A$ \emph{and} with those in $B$: it belongs to the centre of both. The algebras are
\begin{equation}
\A_A=\bigoplus_{\{e_\ell\}}\A_A^{(e)},\qquad\text{centre}=\Linf(\{e_\ell\})\ \text{(boundary electric fluxes)}.
\end{equation}
A state decomposes over flux sectors with probabilities $p_e$. Then the reduced density matrix is block diagonal and
\begin{equation}
S_A=-\sum_ep_e\log p_e+\sum_ep_eS_A^{(e)}.
\end{equation}
The first term is a \emph{classical Shannon entropy} of the boundary flux (the centre), and the second is the average of the entropy within each sector. In the continuum the first term is an additional contribution from the edge modes, the second is the entropy of the dynamical gauge field.

\begin{keyidea}
Gauss's law makes boundary electric fluxes central in the algebra of a region. The entropy then has a ``classical'' piece (Shannon entropy of the centre) in addition to the quantum one, and the Hilbert space is not a tensor product, but a direct sum over fluxes.
\end{keyidea}

\section{Extended Hilbert space and edge modes}

Two ways to cope with the failure of factorization (the literature on this includes work by Buividovic and Polikarpov, Casini--Huerta--Rosabal, Donnelly, and others):
\begin{enumerate}
\item \textbf{Algebraic:} define the entropy using the algebra $\A_A$ with its centre, as above. This is the canonical definition: the entropy is the von Neumann entropy of the state on $\A_A$ with the trace on each block and the Shannon term for the centre.
\item \textbf{Extended Hilbert space:} enlarge the physical Hilbert space by letting the boundary links carry independent edge degrees of freedom so that $\HH_{\rm ext}=\HH_A\otimes\HH_B$ is a true tensor product, and define the physical subspace by gluing. The entropy then agrees with the algebraic one when the gluing is the one forced by Gauss's law.
\end{enumerate}
In the continuum these edge modes have a gravitational counterpart: boundary degrees of freedom associated with the diffeomorphisms that move a cut surface. The quantity $A/4G$ in generalized entropy can be partially understood as such an edge contribution, but this identification is subtle and is not needed for the crossed-product treatment, where the ``edge'' is replaced by the observer's clock.

\section{What this teaches for gravity}

\begin{center}
\renewcommand{\arraystretch}{1.3}
\begin{tabular}{@{}p{4cm}p{4cm}p{4.2cm}@{}}
\toprule
 & Gauge theory & Gravity (with observer)\\
\midrule
Constraint & Gauss law & Hamiltonian/diffeo constraint\\
What commutes with the constraint & electric flux on $\partial A$ & observer dressing (clock energy)\\
Algebra of region & $\A_A$ with non-trivial centre & crossed product $\M\rtimes\R$\\
Entropy & Shannon of centre $+$ average & $\Sgen=A/4G+S_{\rm bulk}$\\
Trace appears because & sum over sectors with weights & group average with $\ee^{x}$\\
\bottomrule
\end{tabular}
\end{center}

\begin{whatwrong}
The analogy is structural, not literal: in gauge theory the algebra of the region is type I (on the lattice) or type III in the continuum with an additional centre. In gravity the new feature is that the constraint generator is the \emph{modular Hamiltonian itself} (or its geometric avatar), which is what turns type III into type II, a stronger effect than adding a centre.
\end{whatwrong}

\section*{Exercises}
\begin{exercise}
For a $\Z_2$ lattice gauge theory (Ising gauge theory) with a single plaquette split across a cut, enumerate the gauge-invariant operators in $A$ and $B$, identify the central element(s) and compute the entropy of a state of your choice using $S=-\sum p_e\log p_e+\sum p_eS^{(e)}$.
\end{exercise}
\begin{exercise}
Show that the algebra generated by $E_\ell$ ($\ell\in\Sigma$) is central in $\A_A$ but not in the full gauge-invariant algebra of the whole lattice.
\end{exercise}
\begin{exercise}
In what sense is $\Linf(\{e_\ell\})$, the centre, the algebraic counterpart of a superselection sector label (\cref{ch:factors})?
\end{exercise}
